% Solo5: one tender process holds one unikernel guest; the entire
% guest ABI is six hypercalls; the host kernel's existing drivers
% do the real device work. Used in M12-L02.
\documentclass[tikz,border=4pt]{standalone}
\usepackage{tikz}
\input{_m10-styles.texinput}
\begin{document}
\begin{tikzpicture}[m10, x=1mm, y=-1mm]
  % One VM: tender process + unikernel guest
  \draw[draw=stackedge, dashed, line width=0.8pt, rounded corners=2.5pt]
        (0,0) rectangle (62,34);
  \node[note, anchor=south west, font=\sffamily\small,
        text=stackedge] at (0,-1) {one VM, one host process};
  \node[safebox, minimum width=56mm, minimum height=9mm,
        font=\sffamily\large, anchor=north west] at (3,3)
        {unikernel};
  \node[warnbox, minimum width=56mm, minimum height=9mm,
        font=\sffamily\large, anchor=north west] at (3,22)
        {Solo5 tender};
  \draw[{Stealth}-{Stealth}, line width=0.9pt, draw=labelgrey]
        (10,12.5) -- (10,21.5);
  % Host kernel and hardware
  \draw[rounded corners=2.5pt, fill=stackfill, draw=stackedge,
        line width=0.7pt] (0,38) rectangle (62,50);
  \node[note, font=\sffamily\large\bfseries, text=stackedge]
        at (31,44) {Linux kernel + KVM};
  \draw[rounded corners=2.5pt, fill=labelgrey!12, draw=labelgrey!70,
        line width=0.7pt] (0,54) rectangle (62,66);
  \node[note, font=\sffamily\large\bfseries, text=labelgrey]
        at (31,60) {Hardware};
  % Right-hand annotations
  \node[note, anchor=west, font=\sffamily\large, text=labelgrey]
        at (66,7.5) {ships no hardware drivers};
  \node[note, anchor=west, font=\sffamily\large, text=labelgrey,
        align=left] at (66,17)
        {six hypercalls:\\console, clock, net, block, yield, exit};
  \node[note, anchor=west, font=\sffamily\large, text=labelgrey]
        at (66,26.5) {loads it, wires its devices, boots it};
  \node[note, anchor=west, font=\sffamily\large, text=labelgrey]
        at (66,44) {existing drivers do the real device work};
\end{tikzpicture}
\end{document}
